% !PRINTED 363
\weberpage{363}
%<<CONT
put. In this none of the three constants $\alpha_1, \beta_2, \gamma_3$ can vanish. For if, e.g., $\alpha_1 = 0$, then by (9) the lines $x_2, x_3, y_1$ intersect on the curve $f$, and this point of intersection would have to be a singular point.

By equation (9), however, the coordinates of the points of contact of $x_1, x_2, x_3$ result from the following three pairs of equations, of which each pair gives two points of contact:
\[
\begin{gathered}
x_1 = 0,\quad x_2(\beta_2 x_2 + \beta_3 x_3) - x_3(\gamma_2 x_2 + \gamma_3 x_3) = 0\\
x_2 = 0,\quad x_3(\gamma_3 x_3 + \gamma_1 x_1) - x_1(\alpha_3 x_3 + \alpha_1 x_1) = 0\\
x_3 = 0,\quad x_1(\alpha_1 x_1 + \alpha_2 x_2) - x_2(\beta_1 x_1 + \beta_2 x_2) = 0.
\end{gathered}
\]
If now these six points are to lie on a conic $\varphi = 0$, then $\varphi$, apart from a constant factor $h$, must pass over, for $x_1 = 0$, into the left part of the second equation of the first pair $x_2(\beta_2 x_2 + \beta_3 x_3) - x_3(\gamma_2 x_2 + \gamma_3 x_3)$. If we therefore denote by $a_1, a_2, a_3$ the coefficients of $x_1^2, x_2^2, x_3^2$ in $\varphi$, then we must have
\weq{10}{
\begin{gathered}
a_2 = h_1 \beta_2,\quad a_3 = -h_1 \gamma_3\\
a_3 = h_2 \gamma_3,\quad a_1 = -h_2 \alpha_1\\
a_1 = h_3 \alpha_1,\quad a_2 = -h_3 \beta_2
\end{gathered}
}
and here $h_1, h_2, h_3$ are three constants.

But from (10) it follows that
\[
h_2 = -h_3,\quad h_3 = -h_1,\quad h_1 = -h_2,
\]
and this would only be possible if $h_1 = h_2 = h_3 = 0$.
But this is the case only when $a_1, a_2, a_3$ vanish, that is, when the conic $\varphi$ passes through the points of intersection of the three straight lines $x_1, x_2, x_3$. But it is to pass through the points of contact of these three double tangents, which are certainly different from their three points of intersection. Thus our assumption is shown to be impossible, and Theorem 3. is proved.

\begin{center}
§. 97.

\gesperrt{Complex pairs and complex triples.}
\end{center}

The theorems of the preceding section give an excellent means of exploring the manifold geometric and algebraic relations of the double tangents and of
%CONT>>
